\section{Conclusion}
\label{sec:conclusion}

This paper began with a conversion and ends with an acknowledgement of
incomplete work.

The conversion was genuine: a practitioner with decades of experience in
distributed systems, skeptical of AI-assisted code generation on structural
grounds, changed his mind in a single morning because the tool contributed
something he had not anticipated. The network monitor worked; the device vendor
lookup was a gift.

What the monitor did not have---and still does not have---is a formal state
model. The system that converted the skeptic remains, as of this writing, a set
of scripts and a Prometheus database. There is no \texttt{.trz} file for
MonitoreoRed. The transition from \texttt{KnownDevice} to
\texttt{Absent} to \texttt{ActiveAlert} is still implicit, still
unverifiable, still invisible to the system that is supposed to track it.

This is not a failure of Trenza. It is the honest boundary of what this paper
can claim. Trenza exists and compiles. It has verified the complete formal
specification of CronometroPSP---thirteen contexts across sixteen \texttt{.trz}
files---in under one hundred milliseconds, detecting completeness violations,
unreachable contexts, and access-control breaches that would otherwise have been
silent runtime defects. It has demonstrated that the presence of a
\texttt{.trz} specification transforms LLM-assisted code review from a
heuristic process into a deductive one. It has been used to specify the very
tool that compiles it.

What remains is the transfer. The language that human-LLM collaboration demanded
into existence must now be applied to the system that motivated its demand. That
application is the second production case, and it is unfinished. We name it
here not as a roadmap commitment but as the open promise that closes the circle
of the argument: the skeptic was converted by a tool that built a system whose
state model the same skeptic later found himself unable to keep in order. The
story does not end with Trenza compiling. It ends---if it ends---when
MonitoreoRed compiles too.

\subsection*{Constraint as gift}

The thesis of this paper is not that Trenza solves the problem of
LLM-generated code. It is that the problem is structural, and that structural
problems require structural solutions. LLMs are not undisciplined because they
lack capability. They are undisciplined because the languages and practices in
which they were trained did not forbid indiscipline.

Trenza imposes constraint on the collaboration---on the human who might scatter
state across four locations, and on the model that would faithfully reproduce
the pattern. That constraint is not a restriction taken away from either party.
It is something both parties receive. A missing handler is easier to fix than
a missing guard condition, precisely because it is visible. A role with a
divergent event surface is easier to reconcile than a role whose divergence is
implicit in scattered code, precisely because the compiler names the divergence
and points at both ends. What the language gives up in expressiveness, it
returns as the ability for both parties to know, at any point, what the other
has committed to.

This is what we mean by \emph{constraint as gift}. The ungiven freedom---the
freedom to write a partial specification, the freedom to leave a behavioral case
unaccounted for, the freedom to let the same name mean two different things in
two different files---is not a freedom either party would choose if the cost of
it were made visible at the time it was exercised. The compiler makes the cost
visible at the time it would be incurred. The collaboration becomes possible
because nothing is implicit between the collaborators.

The adult in the room does not write the code. It makes certain code
unwritable.

This paper was written by the models that needed the adult. The human provided
the room.
